%%%% RESUMO

\begin{resumoutfpr}%% Ambiente resumoutfpr
O Fator Casa representa um dos pilares tradicionais da cultura do futebol, mas
as transformações táticas, logísticas e tecnológicas da era moderna levantam
questionamentos sobre sua real magnitude e sobre os mecanismos que o sustentam.
Este trabalho analisa de forma quantitativa a influência do mando de campo na
Premier League, cobrindo onze temporadas completas, de 2015/2016 a 2025/2026,
em um conjunto de 4.180 partidas. A pesquisa segue a metodologia CRISP-DM e
emprega a linguagem Python para coleta, saneamento, análise estatística e
modelagem. Além da análise descritiva e inferencial, estima-se um modelo de
Poisson bivariado com termo de mando de campo, que permite quantificar a
vantagem do mandante controlando pela força do adversário, e avaliam-se modelos
preditivos sob validação temporal. Os resultados confirmam a existência da
vantagem com alto grau de confiança: mandantes vencem 58,09 por cento das
partidas decididas e obtêm 0,37 ponto a mais por jogo, o que corresponde a um
acréscimo de 21,7 por cento no número esperado de gols. O tamanho de efeito
associado, entretanto, é pequeno, e a vantagem manifesta-se por maior volume de
oportunidades, e não por superioridade na eficiência de finalização. Três
achados contrariam expectativas correntes. Primeiro, a vantagem de mando é
crescente com a qualidade da equipe: clubes que terminam entre os seis primeiros
colocados obtêm diferencial superior ao dos seis últimos, em gradiente
monotônico que se repete em todas as métricas analisadas. Segundo, o viés
disciplinar favorável ao mandante desaparece nas partidas sem público já na
temporada 2019/2020, enquanto a vantagem em pontos permanece inalterada, o que
dissocia o viés de arbitragem do mecanismo central do fenômeno. Terceiro, a
atribuição da queda da vantagem na temporada 2020/2021 à ausência de torcida não
se sustenta quando o contraste é construído dentro de uma mesma temporada. Na
etapa preditiva, o modelo estatístico supera os modelos de aprendizado de
máquina avaliados, e a análise de calibração evidencia por que a acurácia é
insuficiente como critério de avaliação em problemas probabilísticos.
\end{resumoutfpr}
